\section{实直线和实数}

\begin{definition}{实数,实直线,实直线的加法群}{}
	把加法群$\Q$的完备化记作$\R$.
	\begin{enumerate}
		\item $\R$中的元素称为实数.
		\item 拓扑空间$\R$称为实直线,拓扑群$\R$称为实直线的加法群.
		\item 设$\iota:\Q\hookrightarrow\R$是典范映射.我们称$\R\setminus\im\left(\iota\right)$是无理数集,其中的元素称为无理数.
	\end{enumerate}
\end{definition}

\begin{proposition}{$\R$上标准全序的构造}{construction-of-standard-total-ordering-over-R}
	设$\iota:\Q\hookrightarrow\R$是典范映射.定义$\R$上的全序关系$\leq$为
	\begin{align*}
		\forall x,y\in\R\left(x\leq y\iff y-x\in\closure\left(\iota\left(\Q_{>0}\right)\right)\right).
	\end{align*}
	那么,$\leq$是使$\R$构成有序群的全序,$\iota$是保序同构.
\end{proposition}

\begin{definition}{$\R$上的标准全序}{}
	沿用\cref{prop:construction-of-standard-total-ordering-over-R}的设定.我们称$\leq$是$\R$上的全序.
\end{definition}

\begin{convention}{}{}
	无歧义时,我们总是把$\im\left(\Q\right)$视同$\Q$,在此意义下$\Q$是$\R$的子集,$\R\setminus\im\left(\iota\right)$相应地写作$\R\setminus\Q$.
\end{convention}

\begin{proposition}{}{}
	$\R_{\geq 0}$是$\R$中的闭集,$\R_{>0}$是$\R$中的开集.
\end{proposition}